% !TeX root = ../main.tex

\section{Introdução}
\label{sec:introducao}

Nosso trabalho tem como objetivo revisar os conceitos de Pesquisa Operacional aprendidos em sala e sua aplicação em projetos de grande porte envolvendo arquitetura e engenharia. Primeiramente, devemos contextualizar que o desenvolvimento de um empreendimento de arquitetura e engenharia encadeia emissões de projeto, análises e revisões entre disciplinas, cliente e equipes de coordenação. Além disso, as relações entre essas atividades podem gerar atrasos, enquanto a disponibilidade das equipes limita quantas frentes podem avançar simultaneamente. Coordenar as precedências e os recursos é, portanto, parte central do planejamento do prazo. Esse tipo de problema pertence à família do sequenciamento de projetos com restrições de recursos (RCPSP, do inglês *resource-constrained project scheduling problem*), que busca determinar um cronograma que minimize o prazo de conclusão sem violar as dependências nem as capacidades disponíveis \cite{ding2023extensions}.

A questão central da pesquisa é determinar, ao fixar atividades, durações, equipes e vínculos da base, o quanto a escolha dos dias de início pode afetar a diminuição do prazo. Essa comparação parte de uma referência normalizada de 310 dias úteis e distingue duas condições para as atividades sem predecessoras, sendo as largadas livres e datas de liberação preservadas como limites mínimos para os inícios. Esta segunda condição incorpora ao modelo as datas observadas para representar esperas associadas a predecessoras ausentes do recorte, ou seja, não valida essas datas como previsão do prazo real. A primeira condição mostra o limite inferior imposto pelas precedências na base sem liberações e o estudo contribui com uma formulação matemática e uma heurística construtiva, além da verificação independente dos cronogramas principais. A análise refaz o planejamento partindo do inicio considerando, inclusive tarefas já executadas, e não estima o prazo civil do empreendimento nem reprograma sua execução em curso.

O modelo inteiro é útil para formalizar as decisões e produzir limites, entretanto, na instância completa, a busca do CBC terminou após cerca de 36 minutos com uma solução de 307 dias e sem garantir a otimalidade. Tendo em vista que esse resultado está distante do limite inferior de 187 dias, podemos afirmar que ele não respondia com precisão suficiente à pergunta sobre o prazo possível. Por isso, o estudo complementa a formulação com métodos de caminho crítico e uma heurística construtiva do RCPSP. A heurística produziu cronogramas viáveis que atingem o limite inferior nas duas condições de liberação, permitindo provar os ótimos de 285 e 187 dias, respectivamente. A motivação e o papel complementar desses métodos são detalhados nas Seções~\ref{sec:literatura} e~\ref{sec:algoritmo}.

As Seções~\ref{sec:literatura} a~\ref{sec:conclusoes} apresentam os trabalhos relacionados, o escopo e a formulação do problema, os métodos de solução utilizados, os resultados nas duas condições de liberação e as limitações com suas respectivas conclusões.
